\documentclass[11pt,a4paper]{article}
% ---- Préambule commun ----
\usepackage[utf8]{inputenc}
\usepackage[T1]{fontenc}
\usepackage{lmodern}
\usepackage[french]{babel}
\usepackage[a4paper,margin=2.2cm,top=2.6cm,bottom=2.4cm]{geometry}
\usepackage{amsmath,amssymb,amsthm}
\usepackage{enumitem}
\usepackage{xcolor}
\usepackage{listings}
\usepackage{fancyhdr}
\usepackage[most]{tcolorbox}
\usepackage{hyperref}

\definecolor{teal}{HTML}{026573}
\definecolor{tealclair}{HTML}{E6F2F3}
\definecolor{grisclair}{HTML}{F5F5F5}
\definecolor{orangefonce}{HTML}{B35900}

\hypersetup{colorlinks=true,linkcolor=teal,urlcolor=teal,citecolor=teal}

\setlength{\parindent}{0pt}
\setlength{\parskip}{3pt}

% ---- Style Python ----
\lstdefinestyle{python}{
  language=Python,
  basicstyle=\ttfamily\small,
  keywordstyle=\color{teal}\bfseries,
  commentstyle=\color{gray}\itshape,
  stringstyle=\color{orangefonce},
  showstringspaces=false,
  numbers=left,
  numberstyle=\tiny\color{gray},
  numbersep=8pt,
  frame=leftline,
  framerule=1.2pt,
  rulecolor=\color{teal},
  xleftmargin=14pt,
  backgroundcolor=\color{grisclair},
  breaklines=true,
  literate={é}{{\'e}}1 {è}{{\`e}}1 {ê}{{\^e}}1 {à}{{\`a}}1 {ô}{{\^o}}1 {û}{{\^u}}1 {î}{{\^i}}1 {ç}{{\c c}}1
}
\lstset{style=python}

% ---- Compteur d'exercices ----
\newcounter{exo}
\newcommand{\exo}[1]{%
  \stepcounter{exo}%
  \vspace{10pt}%
  \noindent\begin{tcolorbox}[colback=tealclair,colframe=teal,boxrule=1pt,
      left=8pt,right=8pt,top=5pt,bottom=5pt,arc=2pt]
    \textbf{\large Exercice \theexo} \;---\; \textbf{#1}
  \end{tcolorbox}%
  \vspace{4pt}%
  \nopagebreak
}

\newcommand{\partiecours}[1]{%
  \vspace{14pt}%
  {\color{teal}\hrule height 2pt}%
  \vspace{4pt}%
  {\Large\bfseries\color{teal} #1}%
  \vspace{2pt}%
  {\color{teal}\hrule height 0.6pt}%
  \vspace{8pt}%
}

\newtcolorbox{rappel}[1][Rappel]{
  colback=white,colframe=gray!60,boxrule=0.8pt,
  fonttitle=\bfseries,title=#1,coltitle=black,colbacktitle=grisclair,
  left=6pt,right=6pt}

\newtcolorbox{indication}{
  colback=white,colframe=orangefonce!60,boxrule=0.8pt,
  left=6pt,right=6pt,top=4pt,bottom=4pt}

\newcommand{\reponse}{\vspace{4pt}\noindent{\color{teal}\textbf{Solution.}}\;}

\setlength{\headheight}{14pt}
\pagestyle{fancy}
\fancyhf{}
\fancyfoot[C]{\thepage}
\renewcommand{\headrulewidth}{0.4pt}
\renewcommand{\footrulewidth}{0pt}

\fancyhead[L]{\small\color{teal} TD --- Complexité algorithmique}
\fancyhead[R]{\small\color{teal} Énoncé}

\begin{document}

\begin{center}
  {\color{teal}\hrule height 2.5pt}
  \vspace{10pt}
  {\Huge\bfseries\color{teal} Travaux dirigés}\\[6pt]
  {\LARGE\bfseries Calcul de complexité algorithmique}\\[10pt]
  {\large Cas itératif \;$\bullet$\; Cas récursif}\\[8pt]
  {\normalsize CPGE --- filières MP / PSI \quad$|$\quad Informatique pour tous}\\[6pt]
  \vspace{6pt}
  {\color{teal}\hrule height 1pt}
\end{center}

\vspace{6pt}

\begin{rappel}[Conventions et notations]
\begin{itemize}[leftmargin=*,itemsep=2pt]
  \item $n$ désigne toujours la \emph{taille de l'entrée}. On compte le nombre
        d'\emph{opérations élémentaires} (affectations, comparaisons, opérations
        arithmétiques sur des entiers machine), chacune de coût $\Theta(1)$.
  \item $f = O(g)$ : $\exists\, C>0,\ \exists\, n_0,\ \forall n\ge n_0,\ |f(n)|\le C\,g(n)$.
  \item $f = \Omega(g)$ : $g = O(f)$. \qquad $f = \Theta(g)$ : $f=O(g)$ \emph{et} $f=\Omega(g)$.
  \item $f = o(g)$ : $f(n)/g(n)\longrightarrow 0$.
  \item Sommes usuelles : $\displaystyle\sum_{k=1}^{n}k=\frac{n(n+1)}{2}$, \quad
        $\displaystyle\sum_{k=1}^{n}k^2=\frac{n(n+1)(2n+1)}{6}$, \quad
        $\displaystyle\sum_{k=0}^{n}x^k=\frac{x^{n+1}-1}{x-1}\ (x\ne1)$,
  \item Série harmonique : $\displaystyle H_n=\sum_{k=1}^{n}\frac1k=\ln n+\gamma+o(1)$.
  \item Sauf mention contraire, $\log$ désigne le logarithme en base $2$.
\end{itemize}
\end{rappel}

\begin{indication}
\textbf{Méthode de travail.} Pour chaque exercice on demande, dans cet ordre :
(1) l'\emph{expression exacte} du nombre d'opérations lorsqu'elle est accessible ;
(2) un \emph{équivalent} ou un encadrement ; (3) la classe $\Theta$ finale.
Répondre « c'est en $O(n^2)$ » sans justification ne vaut aucun point.
\end{indication}

% =====================================================================
\partiecours{Partie I --- Le cas itératif : boucles \texttt{for} imbriquées}
% =====================================================================

\exo{Échauffement --- compter avant de majorer}
Pour chacun des fragments suivants, donner le nombre \emph{exact} d'exécutions de
l'instruction \texttt{s += 1}, puis la complexité en $\Theta$.

\begin{lstlisting}
# (a)
s = 0
for i in range(n):
    s += 1

# (b)
s = 0
for i in range(n):
    for j in range(n):
        s += 1

# (c)  m et n sont deux entiers independants
s = 0
for i in range(n):
    for j in range(m):
        s += 1

# (d)
s = 0
for i in range(n):
    for j in range(5):
        for k in range(n):
            s += 1
\end{lstlisting}

\medskip
\textbf{Question de synthèse.} Dans (c), peut-on dire que la complexité est
$\Theta(n^2)$ ? Discuter selon la relation entre $m$ et $n$.

% ---------------------------------------------------------------------
\exo{Boucles triangulaires}
Même travail. On donnera à chaque fois la somme exacte \emph{sous forme close}.

\begin{lstlisting}
# (a)
for i in range(n):
    for j in range(i):
        s += 1

# (b)
for i in range(n):
    for j in range(i, n):
        s += 1

# (c)
for i in range(1, n+1):
    for j in range(i, n+1):
        for k in range(j, n+1):
            s += 1
\end{lstlisting}

\begin{enumerate}[label=\arabic*.,leftmargin=*]
  \item Traiter (a) et (b) et vérifier que les deux sommes diffèrent de $n$.
  \item Pour (c), montrer que le nombre d'itérations est le nombre de triplets
        $1\le i\le j\le k\le n$, c'est-à-dire $\binom{n+2}{3}$. En déduire un équivalent.
  \item \textbf{Généralisation.} Quelle est la complexité de $p$ boucles imbriquées
        de ce type ($1\le i_1\le i_2\le\dots\le i_p\le n$) ?
\end{enumerate}

% ---------------------------------------------------------------------
\exo{Indices strictement croissants}
Soit \texttt{T} un tableau de $n$ entiers.

\begin{lstlisting}
# (a) recherche d'un triplet de somme nulle
for i in range(n):
    for j in range(i+1, n):
        for k in range(j+1, n):
            if T[i] + T[j] + T[k] == 0:
                return True
return False

# (b)
for i in range(n):
    for j in range(i+1, n):
        for k in range(n):
            s += 1
\end{lstlisting}

\begin{enumerate}[label=\arabic*.,leftmargin=*]
  \item Pour (a), donner le nombre exact de tests effectués \emph{dans le pire des cas}
        et la complexité correspondante.
  \item Quelle est la complexité de (a) dans le \emph{meilleur} des cas ? Conclusion
        sur l'écriture $\Theta(\cdot)$ pour cet algorithme.
  \item Pour (b), donner le nombre exact d'itérations et la classe $\Theta$.
        Comparer avec (a) : l'ordre de grandeur change-t-il ?
\end{enumerate}

% ---------------------------------------------------------------------
\exo{Bornes non linéaires et pas variable}
Donner pour chaque fragment le nombre d'itérations sous forme de somme, puis un
équivalent et la classe $\Theta$.

\begin{lstlisting}
# (a)
for i in range(n):
    for j in range(i*i):
        s += 1

# (b)
for i in range(1, n+1):
    for j in range(1, n+1, i):     # pas egal a i
        s += 1

# (c)
for i in range(1, n+1):
    j = 0
    while j*j <= i:
        j += 1
        s += 1
\end{lstlisting}

\begin{indication}
Pour (b), le nombre d'itérations internes est $\big\lceil n/i\big\rceil$.
On pensera à la série harmonique. Pour (c), encadrer $\sum_{i=1}^{n}\sqrt{i}$
par des intégrales.
\end{indication}

% ---------------------------------------------------------------------
\exo{Progression géométrique : attention aux fausses évidences}
Les deux fragments ci-dessous ont une boucle externe qui ne fait que
$\Theta(\log n)$ tours. Pourtant leurs complexités diffèrent radicalement.

\begin{lstlisting}
# (a)
i = 1
while i < n:
    for j in range(i):
        s += 1
    i = 2*i

# (b)
i = 1
while i < n:
    for j in range(n):
        s += 1
    i = 2*i
\end{lstlisting}

\begin{enumerate}[label=\arabic*.,leftmargin=*]
  \item Calculer exactement le nombre d'exécutions de \texttt{s += 1} dans (a),
        en supposant $n=2^p$. Conclure. Le résultat vous surprend-il ?
  \item Même question pour (b).
  \item \textbf{Moralité.} Énoncer en une phrase la règle qui distingue les deux cas
        (« le coût total d'une boucle est la \emph{somme} des coûts des tours, pas le
        produit \emph{nombre de tours} $\times$ \emph{coût du dernier tour} »).
  \item Modifier (a) en remplaçant \texttt{i = 2*i} par \texttt{i = i*i} (avec
        \texttt{i} initialisé à $2$). Quelle est alors la complexité ?
\end{enumerate}

% =====================================================================
\partiecours{Partie II --- Le cas itératif : boucles \texttt{while}}
% =====================================================================

\begin{rappel}[Méthode pour une boucle \texttt{while}]
Une boucle \texttt{while} n'a pas de nombre de tours écrit dans le code. On procède
en trois temps :
\textbf{(i)} on identifie un \emph{variant} de boucle, c'est-à-dire une quantité
entière positive strictement décroissante à chaque tour (cela prouve la terminaison) ;
\textbf{(ii)} on exprime la valeur de la variable de contrôle \emph{après $k$ tours} ;
\textbf{(iii)} on résout l'inéquation d'arrêt en $k$.
\end{rappel}

\exo{Les trois schémas fondamentaux}
Pour chaque boucle, exprimer la valeur de la variable après $k$ tours, en déduire le
nombre exact de tours, puis la complexité en fonction de $n$ ($n\ge1$ entier).

\begin{lstlisting}
# (a)
x = n
while x > 0:
    x = x - 1

# (b)  c > 0 constante entiere
x = n
while x > 0:
    x = x - c

# (c)
x = n
while x > 1:
    x = x // 2

# (d)
x = 1
while x < n:
    x = 3*x
\end{lstlisting}

\begin{enumerate}[label=\arabic*.,leftmargin=*]
  \item Traiter (a)--(d). Pour (c), montrer que le nombre de tours vaut exactement
        $\lfloor \log_2 n\rfloor$.
  \item Dans (b), la complexité dépend-elle de $c$ ? Que dit-on alors de $c$ vis-à-vis
        du $\Theta$ ?
  \item \textbf{Piège classique.} Dans (c), l'entrée est l'entier $n$. Si l'on convient
        que la taille de l'entrée est le nombre de \emph{bits} de $n$, soit
        $t=\lfloor\log_2 n\rfloor+1$, quelle est la complexité en fonction de $t$ ?
        Commenter (a) à la lumière de cette convention.
\end{enumerate}

% ---------------------------------------------------------------------
\exo{Algorithme d'Euclide}
\begin{lstlisting}
def pgcd(a, b):
    while b != 0:
        a, b = b, a % b
    return a
\end{lstlisting}

\begin{enumerate}[label=\arabic*.,leftmargin=*]
  \item Montrer que la suite des seconds arguments est strictement décroissante et
        positive : en déduire la terminaison (préciser le variant).
  \item Montrer que si $a > b > 0$ alors $a \bmod b < a/2$.
        \emph{Indication :} distinguer $b\le a/2$ et $b > a/2$.
  \item En déduire qu'après deux tours de boucle, le premier argument a été au moins
        divisé par $2$, puis que le nombre de tours est $O(\log \min(a,b))$.
  \item \textbf{Pire des cas.} On note $F_k$ la suite de Fibonacci ($F_0=0$, $F_1=1$).
        Montrer que l'appel \texttt{pgcd(}$F_{k+1},F_k$\texttt{)} effectue exactement
        $k$ tours de boucle. En déduire que la borne précédente est atteinte et
        exprimer le nombre de tours en fonction de $\varphi=\frac{1+\sqrt5}{2}$.
\end{enumerate}

% ---------------------------------------------------------------------
\exo{Exponentiation rapide}
\begin{lstlisting}
def puissance(x, n):
    r = 1
    while n > 0:
        if n % 2 == 1:
            r = r * x
        x = x * x
        n = n // 2
    return r
\end{lstlisting}

\begin{enumerate}[label=\arabic*.,leftmargin=*]
  \item Prouver la correction en exhibant l'invariant de boucle
        $r\cdot x^{\,n} = x_0^{\,n_0}$ (où $x_0,n_0$ sont les valeurs initiales).
  \item Donner le nombre exact de tours de boucle pour $n\ge1$.
  \item Soit $\nu(n)$ le nombre de $1$ dans l'écriture binaire de $n$. Montrer que le
        nombre total de multiplications vaut $\lfloor\log_2 n\rfloor + 1 + \nu(n)$.
        (On distinguera les élévations au carré et les multiplications accumulées dans
        \texttt{r}.) Quelle multiplication inutile pourrait-on économiser ?
  \item En déduire un encadrement de la complexité, puis la classe $\Theta$ dans le
        pire des cas et dans le meilleur des cas (on précisera pour quels $n$ ils sont
        atteints).
\end{enumerate}

% ---------------------------------------------------------------------
\exo{Deux indices qui se croisent}
Soit \texttt{T} un tableau \emph{trié par ordre croissant} de $n$ entiers deux à deux
distincts et $s$ un entier. L'algorithme suivant cherche deux indices $g<d$ tels que
\texttt{T[g] + T[d] == s}.

\begin{lstlisting}
def paire_de_somme(T, s):
    g, d = 0, len(T) - 1
    while g < d:
        total = T[g] + T[d]
        if total == s:
            return (g, d)
        elif total < s:
            g = g + 1
        else:
            d = d - 1
    return None
\end{lstlisting}

\begin{enumerate}[label=\arabic*.,leftmargin=*]
  \item Proposer un variant de boucle et prouver la terminaison.
  \item Montrer que le nombre de tours est au plus $n-1$, et que cette borne est atteinte.
  \item En déduire la complexité dans le pire des cas. Comparer avec l'algorithme naïf
        à deux boucles \texttt{for} imbriquées.
  \item \emph{(Correction, plus difficile.)} Montrer l'invariant : « si une solution
        $(i,j)$ existe avec $i<j$, alors $g\le i$ et $j\le d$ ».
\end{enumerate}

% =====================================================================
\partiecours{Partie III --- Le cas récursif : équation caractéristique}
% =====================================================================

\begin{rappel}[Récurrences linéaires à coefficients constants]
Soit $(u_n)$ vérifiant, pour $n\ge p$,
\[
  u_n = a_1 u_{n-1} + a_2 u_{n-2} + \dots + a_p u_{n-p} + f(n).
\]
\textbf{1.} On résout d'abord l'équation \emph{homogène} ($f\equiv0$) : on cherche les
racines de l'\emph{équation caractéristique}
$\;r^{p} = a_1 r^{p-1} + \dots + a_p$.
Si $r_1,\dots,r_m$ sont les racines de multiplicités $\mu_1,\dots,\mu_m$, la solution
générale homogène est
$\;u_n^{(h)} = \sum_{i=1}^{m} P_i(n)\, r_i^{\,n}$, avec $\deg P_i < \mu_i$.
\textbf{2.} On cherche une solution \emph{particulière} de la même forme que $f$
(constante si $f$ constante, polynomiale de degré $d$ si $f$ polynomiale de degré $d$),
en majorant le degré de $1$ à chaque fois que la forme essayée est déjà solution de
l'homogène.
\textbf{3.} On somme, puis on détermine les constantes avec les conditions initiales.
\end{rappel}

\exo{Tours de Hanoï et variantes}
On note $T(n)$ le nombre de déplacements effectués par l'algorithme récursif classique.
\begin{lstlisting}
def hanoi(n, depart, arrivee, inter):
    if n == 0:
        return
    hanoi(n-1, depart, inter, arrivee)
    deplacer(depart, arrivee)          # 1 deplacement
    hanoi(n-1, inter, arrivee, depart)
\end{lstlisting}

\begin{enumerate}[label=\arabic*.,leftmargin=*]
  \item Écrire la relation de récurrence vérifiée par $T$ ainsi que sa condition initiale.
  \item Résoudre par la méthode de l'équation caractéristique (homogène + particulière).
        Vérifier le résultat pour $n=1,2,3$.
  \item Résoudre de même les récurrences suivantes (avec $T(0)=0$) :
        \begin{enumerate}[label=(\alph*)]
          \item $T(n) = T(n-1) + n$
          \item $T(n) = T(n-1) + n^2$
          \item $T(n) = 2T(n-1) + n$
          \item $T(n) = 3T(n-1) - 2T(n-2)$, avec $T(0)=1,\ T(1)=3$.
        \end{enumerate}
  \item Pour chacune, donner la classe $\Theta$.
\end{enumerate}

% ---------------------------------------------------------------------
\exo{Coût de Fibonacci naïf}
\begin{lstlisting}
def fib(n):
    if n <= 1:
        return n
    return fib(n-1) + fib(n-2)
\end{lstlisting}

On note $C(n)$ le nombre total d'\emph{appels} à \texttt{fib} engendrés par
\texttt{fib(n)} (l'appel initial compris).

\begin{enumerate}[label=\arabic*.,leftmargin=*]
  \item Établir la récurrence vérifiée par $C$, avec ses deux conditions initiales.
  \item Résoudre par équation caractéristique. On notera
        $\varphi=\frac{1+\sqrt5}{2}$ et $\psi=\frac{1-\sqrt5}{2}$.
  \item Montrer que $C(n) = 2F_{n+1}-1$ où $(F_k)$ est la suite de Fibonacci.
  \item En déduire $C(n)\sim \dfrac{2}{\sqrt5}\,\varphi^{\,n+1}$ et conclure :
        la complexité de \texttt{fib} est $\Theta(\varphi^{\,n})$.
  \item \textbf{Application numérique.} Si une machine effectue $10^{9}$ appels par
        seconde, estimer le temps de calcul de \texttt{fib(80)}.
  \item Que devient la complexité si l'on mémoïse les appels ? Justifier brièvement.
\end{enumerate}

% ---------------------------------------------------------------------
\exo{Racine double et second membre}
Résoudre complètement, puis donner la classe $\Theta$ :

\begin{enumerate}[label=(\alph*),leftmargin=*]
  \item $T(n) = 4T(n-1) - 4T(n-2)$, \quad $T(0)=1$, $T(1)=6$.
  \item $T(n) = 2T(n-1) - T(n-2) + 2$, \quad $T(0)=0$, $T(1)=1$.
  \item $T(n) = T(n-1) + T(n-2) + 1$, \quad $T(0)=T(1)=1$.
  \item $T(n) = 2T(n-1) + 2^{n}$, \quad $T(0)=1$.
        \emph{(Attention : $2$ est racine de l'équation caractéristique.)}
\end{enumerate}

\begin{indication}
Pour (d), chercher une solution particulière sous la forme $\alpha\,n\,2^n$ et non
$\alpha\,2^n$.
\end{indication}

% =====================================================================
\partiecours{Partie IV --- Le cas récursif : récurrences de partition}
% =====================================================================

\begin{rappel}[Récurrence de partition]
Un algorithme « diviser pour régner » découpe une instance de taille $n$ en $a$
sous-instances de taille $n/b$, avec un coût $f(n)$ de découpage et de recombinaison :
\[
  T(n) = a\,T\!\left(\frac{n}{b}\right) + f(n), \qquad a\ge1,\ b>1.
\]
On suppose systématiquement $n=b^{k}$ pour éviter les parties entières (le résultat
asymptotique reste valable dans le cas général). Deux techniques :
\textbf{déroulement} (on itère la relation jusqu'au cas de base et on somme la
contribution de chaque niveau) ou \textbf{arbre des appels} (niveau $i$ :
$a^{i}$ nœuds de coût $f(n/b^{i})$ chacun).
\end{rappel}

\exo{Déroulement à la main --- trois schémas de base}
Aucun théorème général n'est autorisé dans cet exercice : on déroulera la récurrence.
On prendra $n=2^k$ et $T(1)=1$.

\begin{enumerate}[label=\arabic*.,leftmargin=*]
  \item \textbf{Recherche dichotomique.} $T(n) = T(n/2) + 1$.
        Écrire les $k+1$ niveaux, sommer, conclure.
  \item \textbf{Somme des éléments d'un tableau par dichotomie.}
\begin{lstlisting}
def somme(T, g, d):          # somme de T[g..d]
    if g == d:
        return T[g]
    m = (g + d) // 2
    return somme(T, g, m) + somme(T, m+1, d)
\end{lstlisting}
        Établir $T(n)=2T(n/2)+\Theta(1)$ et la résoudre par déroulement. Conclure.
        Cette version est-elle meilleure que la boucle \texttt{for} naïve ?
  \item \textbf{Coût de recombinaison linéaire.} Résoudre $T(n) = 2T(n/2) + n$ par
        déroulement : calculer le coût de chaque niveau, le nombre de niveaux, et conclure.
  \item \textbf{Minimum et maximum simultanés.} On cherche le couple $(\min,\max)$ d'un
        tableau de $n=2^k$ éléments en découpant en deux moitiés. Écrire la récurrence
        du nombre de \emph{comparaisons} $C(n)$ avec $C(2)=1$, et montrer que
        $C(n) = \dfrac{3n}{2} - 2$. Comparer aux $2n-3$ comparaisons de l'algorithme
        itératif naïf.
\end{enumerate}

% ---------------------------------------------------------------------
\exo{Théorème général (Master Theorem)}
\begin{rappel}[Théorème général]
Soient $a\ge1$, $b>1$ et $f$ une fonction positive. Soit $T(n)=aT(n/b)+f(n)$.
On pose $\alpha = \log_b a$. Alors :
\begin{itemize}[leftmargin=*,itemsep=1pt]
  \item \textbf{Cas 1} : si $f(n)=O\!\left(n^{\alpha-\varepsilon}\right)$ pour un
        $\varepsilon>0$, alors $T(n)=\Theta\!\left(n^{\alpha}\right)$.
  \item \textbf{Cas 2} : si $f(n)=\Theta\!\left(n^{\alpha}\log^{c} n\right)$ avec $c\ge0$,
        alors $T(n)=\Theta\!\left(n^{\alpha}\log^{c+1} n\right)$.
  \item \textbf{Cas 3} : si $f(n)=\Omega\!\left(n^{\alpha+\varepsilon}\right)$ pour un
        $\varepsilon>0$ \emph{et} s'il existe $c<1$ tel que $a f(n/b)\le c\,f(n)$ pour
        $n$ assez grand (condition de régularité), alors $T(n)=\Theta\!\left(f(n)\right)$.
\end{itemize}
\end{rappel}

\begin{enumerate}[label=\arabic*.,leftmargin=*]
  \item Appliquer le théorème général à chacune des récurrences suivantes, en
        précisant à chaque fois $a$, $b$, $\alpha=\log_b a$ et le cas invoqué :
        \begin{enumerate}[label=(\alph*),itemsep=1pt]
          \item $T(n) = 9\,T(n/3) + n$
          \item $T(n) = T(2n/3) + 1$
          \item $T(n) = 4\,T(n/2) + n^2$
          \item $T(n) = 3\,T(n/4) + n\log n$
          \item $T(n) = 2\,T(n/2) + n/\log n$
        \end{enumerate}
  \item Pour celle(s) des récurrences ci-dessus qui \emph{échappe(nt)} au théorème,
        expliquer précisément pourquoi aucun des trois cas ne s'applique.
  \item Résoudre $T(n) = 2T(n/2) + n\log n$ par la méthode de l'arbre des appels
        (le théorème ne s'applique pas directement sous la forme du cas 1 ou 3) et
        montrer que $T(n) = \Theta(n\log^2 n)$.
  \item Pourquoi la condition de régularité du cas 3 est-elle nécessaire ?
        Donner l'idée sur l'exemple $f(n)=n^{2}\sin^{2} n$ avec $a=b=2$.
\end{enumerate}

% ---------------------------------------------------------------------
\exo{Multiplications rapides : gagner un appel récursif}
\begin{enumerate}[label=\arabic*.,leftmargin=*]
  \item \textbf{Produit de deux entiers (Karatsuba).} Soient $x$ et $y$ deux entiers de
        $n$ chiffres, $n=2m$. On écrit $x = x_1 10^{m} + x_0$ et $y = y_1 10^{m}+y_0$.
        \begin{enumerate}[label=(\alph*)]
          \item Écrire $xy$ en fonction de $x_1y_1$, $x_1y_0$, $x_0y_1$, $x_0y_0$.
                En déduire la récurrence de l'algorithme naïf et sa complexité.
          \item Vérifier l'identité
                $x_1y_0 + x_0y_1 = (x_1+x_0)(y_1+y_0) - x_1y_1 - x_0y_0$.
                Combien de multiplications récursives suffisent alors ?
          \item En déduire $T(n) = 3\,T(n/2) + \Theta(n)$ et résoudre. Donner la valeur
                numérique de l'exposant $\log_2 3$ à $10^{-3}$ près.
        \end{enumerate}
  \item \textbf{Produit de matrices (Strassen).} L'algorithme de Strassen multiplie deux
        matrices $n\times n$ en $7$ produits de matrices $n/2 \times n/2$ et
        $\Theta(n^2)$ additions.
        \begin{enumerate}[label=(\alph*)]
          \item Écrire la récurrence de l'algorithme par blocs naïf ($8$ produits) et
                la résoudre : retrouve-t-on la complexité de l'algorithme usuel ?
          \item Résoudre la récurrence de Strassen et comparer.
          \item À partir de quelle valeur de $a$ la récurrence $T(n)=a\,T(n/2)+\Theta(n^2)$
                cesse-t-elle d'être en $\Theta(n^2\log n)$ ?
        \end{enumerate}
\end{enumerate}

% ---------------------------------------------------------------------
\exo{Changement de variable}
Lorsque l'argument récursif n'est pas de la forme $n/b$, on se ramène au cas standard
par un changement de variable.

\begin{enumerate}[label=\arabic*.,leftmargin=*]
  \item Résoudre $T(n) = T(\sqrt{n}) + 1$, $T(2)=1$, en posant $n = 2^{k}$ puis
        $S(k)=T(2^{k})$. Conclure : $T(n)=\Theta(\log\log n)$.
  \item Résoudre $T(n) = 2\,T(\sqrt{n}) + \log_2 n$, $T(2)=1$, par la même méthode.
  \item Résoudre $T(n) = \sqrt{n}\; T(\sqrt{n}) + n$ (on posera $U(k)=S(k)/2^{k}$
        après le changement de variable).
  \item \textbf{Application.} Un algorithme sur un ensemble d'entiers de $\{1,\dots,n\}$
        se ramène récursivement à un sous-problème sur $\{1,\dots,\sqrt n\}$ avec un coût
        de recombinaison $\Theta(n)$. Écrire sa récurrence et donner sa complexité.
\end{enumerate}

\vspace{14pt}
{\color{teal}\hrule height 1pt}
\vspace{6pt}
\begin{center}
\small\itshape Fin de l'énoncé --- 16 exercices.
\end{center}

\end{document}
